\section{Présentation des résultats : forme française ou américaine}\label{sec:format}

Les nombres écrits par \tkzname{tkz-fn} (graduations des axes, coordonnées, résultats de calculs) sont présentés selon l'une des deux formes suivantes :

\begin{center}
\begin{tabular}{lll}
\toprule
& forme française (\texttt{fr}) & forme américaine (\texttt{us}) \\
\midrule
séparateur décimal & $3{,}14$ & $3.14$ \\
séparateur des milliers & $12\,345$ (espace fine) & $12{,}345$ \\
intervalles & $\mathopen{]}a\,;\,b\mathclose{[}$, $\mathopen{[}a\,;\,b\mathclose{]}$ & $(a,\,b)$, $[a,\,b]$ \\
coordonnées & $A(2\,;\,3)$ & $A(2,\,3)$ \\
origine du repère & $O$ & $0$ \\
\bottomrule
\end{tabular}
\end{center}

Les milliers ne sont séparés qu'à partir de cinq chiffres : $1234$ mais $\tkzFnPrint{12345}$.

\subsection{Choix de la forme}

\begin{itemize}
  \item Par défaut, la forme est française si la langue de \tkzname{babel} est le français au début du document, américaine sinon.
  \item Les options du package imposent la forme : |\usepackage[fr]{tkz-fn}| ou |\usepackage[us]{tkz-fn}|.
  \item Dans le document : |\tkzFnSetup{format=us}|.
\end{itemize}

\begin{NewMacroBox}{tkzFnSetup}{\marg{options}}
\begin{tabular}{lll}
\toprule
options             & défaut & définition  \\
\midrule
\TOline{format}{fr ou us}{forme des nombres}
\TOline{decimals}{2}{nombre maximum de décimales de \tkzcname{tkzFnPrint}}
\TOline{axes}{school}{présentation des axes : \texttt{school} ou \texttt{box}}
\TOline{arrows}{end}{flèches des axes : \texttt{end}, \texttt{both} ou \texttt{none}}
\bottomrule
\end{tabular}
\end{NewMacroBox}

\subsection{Écrire un résultat : \tkzcname{tkzFnPrint}}

\begin{NewMacroBox}{tkzFnPrint}{\oarg{décimales}\marg{expression Lua}}
\emph{Calcule l'expression et écrit le résultat arrondi, sans zéros inutiles. Les fonctions définies par \tkzcname{tkzFnDef} sont utilisables, ainsi que \texttt{diff(f,x)}, \texttt{integral(f,a,b)} et \texttt{riemann(f,a,b,n,"left")}.}
\end{NewMacroBox}

\begin{NewMacroBox}{tkzFnPrintInterval}{\oarg{oo, cc, oc ou co}\marg{a}\marg{b}}
\emph{Écrit un intervalle ; \texttt{o} = ouvert, \texttt{c} = fermé. Les bornes sont des expressions Lua ou du texte \TeX\ (\texttt{-\textbackslash infty}).}
\end{NewMacroBox}

\subsection{Exemple}

\begin{tkzexample}[latex=6cm]
\tkzFnDef{f}{x^2-2*x-1}
\begin{tikzpicture}[scale=.5]
  \tkzFnInit[xmin=-2,xmax=4,
             ymin=-3,ymax=6]
  \tkzFnGrid \tkzFnAxes
  \tkzFnDraw[thick,red](f)
  \tkzFnDefPoint[draw](f,3.5){A}
\end{tikzpicture}

$A\tkzFnPrintCoord(A)$,
$f'(3)=\tkzFnPrint{diff(f,3)}$,
$\tkzFnPrint[4]{f(0.25)}$,
$\tkzFnPrint{1234567.891}$,
$I=\tkzFnPrintInterval[oc]{-\infty}{2.5}$

\tkzFnSetup{format=us}
$A\tkzFnPrintCoord(A)$,
$f'(3)=\tkzFnPrint{diff(f,3)}$,
$\tkzFnPrint[4]{f(0.25)}$,
$\tkzFnPrint{1234567.891}$,
$I=\tkzFnPrintInterval[oc]{-\infty}{2.5}$
\tkzFnSetup{format=fr}
\end{tkzexample}

\subsection{Graduations en forme américaine}

\begin{tkzexample}[latex=7cm]
\tkzFnSetup{format=us}
\begin{tikzpicture}[scale=.8]
  \tkzFnInit[xmin=0,xmax=2,xstep=.25,
             ymin=0,ymax=20000,ystep=4000]
  \tkzFnGrid
  \tkzFnAxes
  \tkzFnDraw[blue,thick]{5000*x^2}
\end{tikzpicture}
\tkzFnSetup{format=fr}
\end{tkzexample}

\endinput
